\documentclass[11pt,a4paper]{article}
\usepackage[utf8]{inputenc}
\usepackage{amsmath,amssymb,amsfonts,amsthm}
\usepackage{geometry}
\geometry{margin=1in}
\usepackage{hyperref}
\usepackage{cite}
\usepackage{bm}

\title{\textbf{A Major-Minor Mean Field Game Heuristics for Swarm Pursuit-Evasion Dynamics}}
\author{\textbf{Mathematical Formulation \& Algorithmic Resolution}}
\date{\today}

\begin{document}

\maketitle

\begin{abstract}
This document establishes a formal mathematical framework for modeling the interaction between a continuum of minor players (a drone swarm) and an active, strategic major player (an evading target) within a bounded spatial domain containing obstacles. The system is formulated as a Major-Minor Pursuit-Evasion Mean Field Game (PE-MFG). We detail the backward Hamilton-Jacobi-Bellman (HJB) equation governing optimal drone navigation, the forward Kolmogorov-Fokker-Planck (KFP) equation modeling swarm mass transport, and the state-costate differential equations determining the evader's optimal escape trajectory. Finally, an extended damped Picard fixed-point scheme is presented for numerical computation.
\end{abstract}

\section{Introduction \& Domain Configuration}
Let $\Omega \subset \mathbb{R}^2$ be a bounded domain with boundary $\partial\Omega = \Gamma_{\text{wall}} \cup \Gamma_{\text{door}}$, where $\Gamma_{\text{wall}}$ represents solid obstacles enforcing zero-flux Neumann conditions and $\Gamma_{\text{door}}(t)$ denotes time-dependent exit/target boundaries. We model a time horizon $t \in [0, T]$.

The system consists of two distinct agent classes:
\begin{enumerate}
    \item \textbf{Minor Players (Drone Swarm):} A continuum of statistically identical agents whose macroscopic spatial density at time $t$ and location $x \in \Omega$ is given by $m(t, x) \ge 0$.
    \item \textbf{Major Player (Evader/Target):} A single dominant agent located at position $y(t) \in \Omega$, actively adjusting its velocity $a(t) \in \mathbb{R}^2$ to escape the approaching swarm.
\end{enumerate}

\section{Minor Player (Swarm) Dynamics}

Each individual drone minimizes a cost functional penalizing running distance to the dynamic target $y(t)$, spatial congestion, and control effort.

\subsection{Hamilton-Jacobi-Bellman (HJB) Equation}
Let $u(t, x)$ denote the value function (expected cost-to-go) of a drone at position $x$ and time $t$. The value function satisfies the backward-in-time HJB equation \cite{caines2017mean}:
\begin{equation}
\label{eq:hjb}
-\frac{\partial u}{\partial t} - \sigma \Delta u + H(x, m, \nabla u) = f(x, y(t)), \quad (t, x) \in (0, T) \times \Omega,
\end{equation}
with terminal condition $u(T, x) = g(x, y(T))$ and Dirichlet boundary condition $u(t, x) = 0$ for $x \in \Gamma_{\text{door}}(t)$.

The congestive Hamiltonian $H(x, m, p)$ incorporates crowd density penalties \cite{carmona2016probabilistic}:
\begin{equation}
H(x, m, p) = \frac{c_0 |p|^2}{(1 + m)^\alpha} - f_0,
\end{equation}
where $\alpha \in (0, 1]$ represents the congestion exponent and $c_0, f_0 > 0$ are physical parameters.

The running cost $f(x, y(t))$ penalizes the spatial separation from the target:
\begin{equation}
f(x, y(t)) = \frac{\chi}{2} \|x - y(t)\|^2,
\end{equation}
where $\chi > 0$ is the attraction weight.

\subsection{Kolmogorov-Fokker-Planck (KFP) Equation}
The aggregate density $m(t, x)$ of the swarm evolves forward in time according to the controlled drift velocity $\mathbf{v}(t, x) = -\nabla_p H(x, m, \nabla u)$ \cite{huang2016backward}:
\begin{equation}
\label{eq:kfp}
\frac{\partial m}{\partial t} - \sigma \Delta m - \text{div}\left( m \, \frac{2 c_0 \nabla u}{(1 + m)^\alpha} \right) = 0, \quad (t, x) \in (0, T) \times \Omega,
\end{equation}
subject to initial density $m(0, x) = m_0(x)$, homogeneous Neumann boundary conditions on $\Gamma_{\text{wall}}$, and Dirichlet absorption ($m(t, x) = 0$) on $\Gamma_{\text{door}}(t)$.

\section{Major Player (Evader) Optimal Control}

The target's trajectory $y(t)$ evolves according to its velocity control $a(t) \in \mathbb{R}^2$:
\begin{equation}
\dot{y}(t) = a(t), \quad y(0) = y_0.
\end{equation}

\subsection{Target Cost Functional}
The target seeks to maximize its distance from high swarm density regions while minimizing kinetic expenditure \cite{bergault2024mean, wan2021improved}:
\begin{equation}
\label{eq:target_cost}
J_E(a) = \int_0^T \left( \int_{\Omega} \Phi(x - y(t)) m(t, x) \, dx + \frac{\gamma}{2} \|a(t)\|^2 \right) dt - G(y(T)),
\end{equation}
where $\Phi(r)$ is a repelling interaction kernel (e.g., Gaussian interaction $\Phi(r) = w_0 \exp(-\|r\|^2 / (2 r_0^2))$) and $\gamma > 0$ denotes control effort cost.

\subsection{Pontryagin's Maximum Principle \& Escape Force}
Defining the target costate vector $p_E(t) \in \mathbb{R}^2$, Pontryagin's Maximum Principle yields the optimal escape velocity \cite{badakaya2022game}:
\begin{equation}
a^*(t) = -\frac{p_E(t)}{\gamma}.
\end{equation}

The costate vector $p_E(t)$ satisfies the adjoint backward differential equation:
\begin{equation}
\label{eq:costate}
\dot{p}_E(t) = \int_{\Omega} \nabla_y \Phi(x - y(t)) m(t, x) \, dx, \quad p_E(T) = -\nabla_y G(y(T)).
\end{equation}

Differentiating the trajectory equation yields a second-order equation of motion driven by the macroscopic threat field:
\begin{equation}
\label{eq:target_acceleration}
\ddot{y}(t) = -\frac{1}{\gamma} \int_{\Omega} \nabla_y \Phi(x - y(t)) m(t, x) \, dx.
\end{equation}

\section{The Coupled Major-Minor Nash Equilibrium}

A Mean Field Nash Equilibrium is defined as a tuple $(u^*, m^*, y^*, p_E^*)$ simultaneously satisfying the coupled system:
\begin{equation}
\begin{cases}
-\dfrac{\partial u}{\partial t} - \sigma \Delta u + H(x, m, \nabla u) = \dfrac{\chi}{2} \|x - y(t)\|^2, & u(T, x) = g(x, y(T)), \\[10pt]
\dfrac{\partial m}{\partial t} - \sigma \Delta m - \text{div}\left( m \nabla_p H(x, m, \nabla u) \right) = 0, & m(0, x) = m_0(x), \\[10pt]
\dot{y}(t) = -\dfrac{p_E(t)}{\gamma}, & y(0) = y_0, \\[10pt]
\dot{p}_E(t) = \int_{\Omega} \nabla_y \Phi(x - y(t)) m(t, x) \, dx, & p_E(T) = -\nabla_y G(y(T)).
\end{cases}
\end{equation}

\section{Algorithmic Resolution via Extended Picard Iteration}

To compute the coupled equilibrium, we employ an iterative fixed-point algorithm with under-relaxation parameter $\theta \in (0, 1]$:

\begin{enumerate}
    \item \textbf{Initialization:} Set initial guesses $U^{(0)}(t, x)$, $M^{(0)}(t, x)$, and $y^{(0)}(t)$ for $t \in [0, T]$.
    \item \textbf{Step 1 (HJB Backward):} Solve Eq.~\eqref{eq:hjb} backward from $t = T$ to $t = 0$ using $M^{(k-1)}$ and $y^{(k-1)}$ to obtain $\tilde{U}^{(k)}$. Update $U^{(k)} = \theta \tilde{U}^{(k)} + (1-\theta) U^{(k-1)}$.
    \item \textbf{Step 2 (KFP Forward):} Solve Eq.~\eqref{eq:kfp} forward from $t = 0$ to $t = T$ using $U^{(k)}$ to obtain $\tilde{M}^{(k)}$. Update $M^{(k)} = \theta \tilde{M}^{(k)} + (1-\theta) M^{(k-1)}$.
    \item \textbf{Step 3 (Target Trajectory Forward/Backward):} Integrate Eq.~\eqref{eq:costate} backward and trajectory $\dot{y}(t)$ forward using $M^{(k)}$ to obtain $\tilde{y}^{(k)}(t)$. Update $y^{(k)} = \theta \tilde{y}^{(k)} + (1-\theta) y^{(k-1)}$.
    \item \textbf{Convergence Check:} Repeat until $\|U^{(k)} - U^{(k-1)}\|_2 < \epsilon$ and $\|M^{(k)} - M^{(k-1)}\|_2 < \epsilon$.
\end{enumerate}

\begin{thebibliography}{99}

\bibitem{caines2017mean}
P.~E. Caines, M.~Huang, and R.~P. Malham{\'e}, 
\emph{Mean field games}, 
Handbook of Dynamic Game Theory, Springer, pp.~345--372, 2017.

\bibitem{carmona2016probabilistic}
R.~Carmona and X.~Zhu, 
\emph{A probabilistic approach to mean field games with major and minor players}, 
The Annals of Applied Probability, vol.~26, no.~3, pp.~1535--1580, 2016.

\bibitem{huang2016backward}
J.~Huang, S.~Wang, and Z.~Wu, 
\emph{Backward-forward linear-quadratic mean-field games with major and minor agents}, 
Probability, Uncertainty and Quantitative Risk, vol.~1, no.~1, p.~8, 2016.

\bibitem{bergault2024mean}
P.~Bergault, P.~Cardaliaguet, and C.~Rainer, 
\emph{Mean field games in a Stackelberg problem with an informed major player}, 
SIAM Journal on Control and Optimization, vol.~62, no.~3, pp.~1737--1765, 2024.

\bibitem{badakaya2022game}
A.~J. Badakaya, A.~S. Halliru, and J.~Adamu, 
\emph{Game value for a pursuit-evasion differential game problem in a Hilbert space}, 
Journal of Dynamics \& Games, vol.~9, no.~1, pp.~1--15, 2022.

\bibitem{wan2021improved}
K.~Wan, D.~Wu, Y.~Zhai, B.~Li, X.~Gao, and Z.~Hu, 
\emph{An improved approach towards multi-agent pursuit--evasion game decision-making using deep reinforcement learning}, 
Entropy, vol.~23, no.~11, p.~1433, 2021.

\end{thebibliography}

\end{document}